\documentclass[11pt]{article}
\usepackage{amsmath,amssymb,amsthm,booktabs,enumitem,geometry}
\geometry{margin=1in}
\newtheorem{theorem}{Theorem}
\newtheorem{lemma}{Lemma}
\newtheorem{proposition}{Proposition}
\newcommand{\C}{\mathbb C}
\newcommand{\eps}{\varepsilon}
\title{Step 125: Prime-Conductor Finite Source-Frame Model}
\author{RATCHET / Six Birds RH Membrane Program}
\date{}
\begin{document}
\maketitle

\section{Purpose}
Step 125 instantiates the conductor-decoupled source theorem on the cleanest finite arithmetic family: complete Dirichlet characters modulo a prime conductor $q_N$ with coefficient support length $L_N<q_N$.

The goal is not to re-prove standard analytic number theory.  The goal is to place the standard finite character orthogonality theorem into the membrane source-frame language and to expose the first defects that matter for the completed residual sector.

\section{Finite coefficient space}
Let $q$ be prime and let
\[
  I_L=\{1,\dots,L\},\qquad L<q.
\]
Let $a=(a_n)_{n\in I_L}\in \C^{I_L}$ and define
\[
  Q_\chi a=\sum_{n\le L} a_n\chi(n).
\]
The full source Gram is
\[
  G_q(m,n)=\sum_{\chi\bmod q}\chi(n)\overline{\chi(m)}.
\]

\begin{lemma}[Complete prime-modulus character frame]
If $q$ is prime and $1\le m,n\le L<q$, then
\[
  G_q(m,n)=(q-1)\delta_{mn}.
\]
Thus
\[
  G_q=(q-1)I_L.
\]
\end{lemma}

\begin{proof}
Characters modulo $q$ are the dual group of $(\mathbb Z/q\mathbb Z)^\times$.  Orthogonality gives
\[
\sum_{\chi\bmod q}\chi(n)\overline{\chi(m)}
=(q-1)\mathbf 1_{n\equiv m\bmod q}.
\]
Since $1\le m,n<q$, congruence modulo $q$ is equality.
\end{proof}

\section{Primitive-character defect}
For prime $q$, all non-principal characters are primitive and the principal character has conductor $1$.  Removing the principal character gives
\[
  G_q^{\rm prim}=G_q-\mathbf 1\mathbf 1^*=(q-1)I_L-J_L.
\]

\begin{proposition}[Primitive rank-one correction]
The eigenvalues of $G_q^{\rm prim}$ on $\C^{I_L}$ are
\[
  q-1-L
\]
for the constant vector, and
\[
  q-1
\]
with multiplicity $L-1$.  Hence
\[
  G_q^{\rm prim}\succeq (q-1-L)I_L.
\]
After averaging over primitive characters,
\[
  \frac{1}{q-2}G_q^{\rm prim}\succeq \left(1-\frac{L}{q-2}\right)I_L.
\]
\end{proposition}

This is the first honest defect: the primitive-only family loses a rank-one constant direction.  Robust primitive lower frame requires $L$ to remain strictly below the conductor scale.

\section{Membrane source-frame implication}
Let $Y_{R,N}$ be a finite residual sector with Gram $G_{R,N}$.  Let
\[
  R_N:Y_{R,N}\to\C^{I_N}
\]
be the regularized Burnol-to-Dirichlet coefficient map.  If
\[
  G_{\mathcal X,N}\succeq \gamma_N H_N
\]
and
\[
  R_N^*H_NR_N\succeq c_{\rm hyb,N}G_{R,N},
\]
then
\[
  R_N^*G_{\mathcal X,N}R_N\succeq \gamma_N c_{\rm hyb,N}G_{R,N}.
\]
Thus the effective source strength is
\[
  \Lambda_{R,N}=\gamma_N c_{\rm hyb,N}.
\]

\section{Source-weight normalization warning}
The unnormalized full prime-modulus frame has $\gamma_N=q_N-1$.  The averaged full frame has $\gamma_N=1$.  Therefore divergence of $\Lambda_N$ is not a free consequence of having many characters.

The P6 audit must declare whether summing over all characters is a lawful accumulation of independent source records, or whether the family must be normalized by total source currency.  Uncharged character-count growth is smuggling.

\section{Completed promotion}
A finite coefficient frame is not a completed membrane.  It must promote through a fixed or exhaustive tail record
\[
  G_R\preceq \Pi_N^*G_{R,N}\Pi_N+T_N.
\]
If the lifted frame is used in the form
\[
  F_N+\Lambda_NT_N\succeq \Lambda_NG_R,
\]
then the amplified tail $\Lambda_NT_N$ must be explicitly controlled.  Otherwise the result remains moving-window support evidence.

\section{Route status}
Step 125 supplies a finite exact source-frame prototype.  The remaining load-bearing work is:
\begin{enumerate}[leftmargin=1.5em]
\item source-weight normalization / cumulative conductor ladder audit;
\item matrix lower-moment estimates for source-weighted families;
\item regularized Burnol-to-Dirichlet visibility floor;
\item fixed/exhaustive residual-tail promotion;
\item all-six and $\Xi$ adequacy records.
\end{enumerate}

\end{document}